% =============================================================================
% Temperature Control Case Study: Methodology and Evaluation
% =============================================================================
% Required packages (add to main thesis preamble if not already present):
%   \usepackage{algorithm}
%   \usepackage{algorithmic}
%   \usepackage{multirow}
% =============================================================================

\chapter{Temperature Control Case Study}
\label{ch:temperature-case-study}

This chapter presents a single-building-cluster temperature control case study
that adapts the STEMS reward framework~\cite{zhang2025stems} to a
safety-constrained single-agent setting.  A Constrained Soft Actor-Critic with
smoothed log-barrier penalty (\textbf{CSAC-LB})~\cite{zhang2024csaclb} is
compared against a bang-bang thermostat baseline (\textbf{RBC}) representative
of commodity HVAC on/off controllers.

% =============================================================================
\section{Methodology}
\label{sec:temp-methodology}

% -----------------------------------------------------------------------------
\subsection{Simulation Environment}
\label{sec:temp-env}

Experiments are conducted on CityLearn v2~\cite{nweye2023citylearn}, a
multi-building energy management simulation platform that provides realistic
building energy dynamics and grid interactions.  The experimental dataset is
derived from real building energy consumption data in Travis County, Texas.  The
dataset spans from August 2018 to August 2019, capturing a full year of
operation under typical subtropical humid climate conditions: annual average
temperature of $20.7\,^\circ$C, temperature range from $-7.8\,^\circ$C to
$42.8\,^\circ$C.

The case-study schema binds three residential buildings with LSTM-enabled indoor
temperature dynamics.  Each building is equipped with a cooling device,
electrical storage, and domestic hot water (DHW) storage.  A centralized
single-agent controller issues actions for all three buildings at hourly
resolution, producing $T = 2{,}207$ active time steps per episode after an
initial 13-step LSTM warm-up period during which comfort costs are disabled.

% -----------------------------------------------------------------------------
\subsection{Problem Formulation}
\label{sec:temp-formulation}

The temperature control problem is formulated as a Constrained Markov Decision
Process (CMDP) $\mathcal{M} = (\mathcal{S}, \mathcal{A}, P, \mathcal{R},
\mathcal{C}, \gamma)$.

\paragraph{State space.}
The observation vector $\mathbf{s}_t \in \mathcal{S} \subset \mathbb{R}^{54}$
is constructed from normalised per-building, environmental, and temporal
features.  Table~\ref{tab:observation-space} details the observation
components.  All features are normalised to $[0, 1]$ by CityLearn's internal
observation wrapper using dataset-derived bounds.

\begin{table}[htbp]
\centering
\caption{Observation space composition ($\dim(\mathcal{S}) = 54$).  Each
  per-building feature is replicated for all $N=3$ buildings.}
\label{tab:observation-space}
\begin{tabular}{llc}
\toprule
\textbf{Category} & \textbf{Feature} & \textbf{Dim} \\
\midrule
\multirow{4}{*}{Building state}
  & Indoor dry-bulb temperature $T_t^{(i)}$ & $3$ \\
  & Temperature setpoint / reference & $3$ \\
  & Cooling device status & $3$ \\
  & Electrical storage state of charge & $3$ \\
  & DHW storage state of charge & $3$ \\
\midrule
\multirow{3}{*}{Environmental}
  & Outdoor dry-bulb temperature & $1$ \\
  & Direct solar irradiance & $1$ \\
  & Diffuse solar irradiance & $1$ \\
  & Solar generation $p_t^{(i)}$ & $3$ \\
  & Electricity price $v_t$ & $1$ \\
\midrule
\multirow{3}{*}{Temporal}
  & Hour of day (sine/cosine encoded) & $2$ \\
  & Day of year (sine/cosine encoded) & $2$ \\
  & Occupancy count per building & $3$ \\
\midrule
\multirow{2}{*}{Grid / district}
  & Net district electricity consumption & $1$ \\
  & District import power & $1$ \\
  & Net building consumption $e_t^{(i)}$ & $3$ \\
\midrule
\multicolumn{2}{l}{\textit{Remaining CityLearn default features}} & $\sim$20 \\
\bottomrule
\multicolumn{2}{l}{\textbf{Total}} & \textbf{54} \\
\end{tabular}
\end{table}

The observation space provides the agent with sufficient information to reason
about thermal dynamics (indoor and outdoor temperatures), energy markets
(electricity price), renewable availability (solar generation), and temporal
patterns (time encoding).  The 13-step LSTM warm-up period at the start of
each episode allows CityLearn's internal LSTM temperature model to stabilise
before the agent's actions begin affecting indoor temperatures.

\paragraph{Action space.}
In the \emph{cooling-only} interface adopted for this case study, the policy
controls a single continuous cooling action per building,
$\mathbf{a}_t \in [0,1]^3$, where $a_t^{(i)} = 0$ means no cooling and
$a_t^{(i)} = 1$ means maximum cooling for building~$i$.
Table~\ref{tab:action-space} summarises the action interface.

\begin{table}[htbp]
\centering
\caption{Action space definition.  Only cooling actions are controlled by the
  RL policy; all other actuators are frozen to isolate the temperature control
  problem.}
\label{tab:action-space}
\begin{tabular}{lccc}
\toprule
\textbf{Actuator} & \textbf{Range} & \textbf{Controlled by} & \textbf{Dim} \\
\midrule
Cooling device & $[0, 1]$ & RL policy & $3$ \\
Electrical storage & $0$ (fixed) & Wrapper & $3$ \\
DHW storage & $0$ (fixed) & Wrapper & $3$ \\
\midrule
\multicolumn{3}{l}{\textbf{Effective policy action dimension}} & \textbf{3} \\
\bottomrule
\end{tabular}
\end{table}

Battery and DHW storage actions are set internally by the environment wrapper
to fixed zero values, ensuring that the CSAC-LB policy operates on a clean
3-dimensional cooling-only action interface.  This isolation is critical for
attributing performance differences to the cooling control strategy rather
than storage management.

\paragraph{Reward function.}
Following the STEMS framework~\cite{zhang2025stems}, the reward is decomposed
into four components:
\begin{equation}
  R_t = R_t^{\text{economic}} + R_t^{\text{stability}} +
        R_t^{\text{renewable}} + R_t^{\text{comfort}}
  \label{eq:reward-decomposition}
\end{equation}

The \emph{economic} component minimises electricity costs:
\begin{equation}
  R_t^{\text{economic}} = -\mu \cdot v_t \cdot \sum_{i=1}^{N} e_t^{(i)}
  \label{eq:reward-economic}
\end{equation}
where $v_t$ is the electricity price, $e_t^{(i)}$ is the net electricity
consumption of building~$i$, and $\mu = 1.0$ is the economic weight.

The \emph{stability} component combines grid-level coordination and
building-level smooth operation:
\begin{equation}
  R_t^{\text{stability}} = \underbrace{N \cdot \alpha_{\text{grid}} \left(
    1 - \left(\frac{\sum_i \max(0, e_t^{(i)})}{P_{\text{grid,max}}}\right)^2
  \right)}_{\text{grid stability}}
  + \underbrace{\alpha_{\text{build}} \sum_{i=1}^{N} \left(
    1 - \frac{|e_t^{(i)}|}{P_{\text{build,max}}}
  \right)}_{\text{building stability}}
  - \underbrace{\beta_{\text{ramp}} \sum_{i=1}^{N} \frac{|e_t^{(i)} - e_{t-1}^{(i)}|}{P_{\text{build,max}}}}_{\text{ramping penalty}}
  \label{eq:reward-stability}
\end{equation}
where $\alpha_{\text{grid}} = 0.5$, $\alpha_{\text{build}} = 0.3$,
$\beta_{\text{ramp}} = 0.2$, and $P_{\text{grid,max}}$, $P_{\text{build,max}}$
are calibrated from the dataset.

The \emph{renewable energy} component maximises clean energy utilisation:
\begin{equation}
  R_t^{\text{renewable}} = \xi \sum_{i=1}^{N}
    \min\!\left(\frac{p_t^{(i)}}{p_t^{(i)} + \max(0, e_t^{(i)})},\; 1\right)
  \label{eq:reward-renewable}
\end{equation}
where $p_t^{(i)}$ is the solar generation and $\xi = 0.6$.

The \emph{comfort} component penalises temperature deviations from the comfort
band $[T_{\min}, T_{\max}] = [20\,^\circ\text{C}, 26\,^\circ\text{C}]$ with
occupancy-aware weighting:
\begin{equation}
  R_t^{\text{comfort}} = -\lambda_{\text{indoor}} \sum_{i=1}^{N}
    w_i \cdot \Delta T_t^{(i)}
  \label{eq:reward-comfort}
\end{equation}
where $\Delta T_t^{(i)} = \max(T_t^{(i)} - T_{\max}, 0) + \max(T_{\min} -
T_t^{(i)}, 0)$ is the comfort deviation, $w_i = 1.0$ if the building is
occupied and $w_i = \eta = 0.1$ otherwise, and $\lambda_{\text{indoor}} = 1.2$
to provide a strong comfort gradient through the barrier penalty.

\paragraph{Cost function.}
The CMDP cost is defined as a weighted district-level comfort violation signal:
\begin{equation}
  C_t = w_{\text{comfort}} \cdot \sum_{i=1}^{N} \mathbf{1}\!\left[
    T_t^{(i)} < T_{\min} \;\text{or}\; T_t^{(i)} > T_{\max}
  \right]
  \label{eq:cmdp-cost}
\end{equation}
where $w_{\text{comfort}} = 0.05$.  During the 13-step LSTM warm-up,
$C_t = 0$.

% -----------------------------------------------------------------------------
\subsection{Background: Soft Actor-Critic (SAC)}
\label{sec:temp-sac-background}

Before introducing the constrained variant, we briefly review the Soft
Actor-Critic (SAC) algorithm~\cite{haarnoja2018sac} that forms the
foundation of our approach.

SAC is an off-policy actor-critic algorithm that augments the standard RL
objective with an entropy bonus, encouraging exploration and robustness:
\begin{equation}
  \pi^* = \arg\max_\pi \sum_{t=0}^{T}
    \mathbb{E}_{(s_t, a_t) \sim \rho_\pi}\!\left[
      \gamma^t \bigl(R(s_t, a_t) + \alpha\, H(\pi(\cdot|s_t))\bigr)
    \right]
  \label{eq:sac-objective}
\end{equation}
where $H(\pi(\cdot|s)) = -\mathbb{E}_{a \sim \pi}[\log \pi(a|s)]$ is the
policy entropy and $\alpha > 0$ is the entropy temperature that controls the
exploration--exploitation trade-off.  The entropy regularisation prevents
premature convergence to deterministic policies, which is particularly
beneficial in building control where the reward landscape contains multiple
local optima corresponding to different cooling strategies.

\paragraph{Critic networks.}
SAC maintains two Q-value networks $Q_{\theta_1}$ and $Q_{\theta_2}$ (the
\emph{twin critics}) trained to minimise the soft Bellman residual:
\begin{equation}
  \mathcal{L}_Q(\theta_j) = \mathbb{E}_{(s,a,r,s') \sim \mathcal{D}}\!\left[
    \bigl(Q_{\theta_j}(s,a) - y\bigr)^2
  \right], \quad j \in \{1, 2\}
  \label{eq:sac-critic-loss}
\end{equation}
where $\mathcal{D}$ is the replay buffer and the target value includes the
entropy bonus:
\begin{equation}
  y = r + \gamma(1 - d_t)\!\left[
    \min_{j=1,2} Q_{\bar{\theta}_j}(s', a')
    - \alpha \log \pi_\phi(a'|s')
  \right], \quad a' \sim \pi_\phi(\cdot|s')
  \label{eq:sac-target}
\end{equation}
Here $\bar{\theta}$ denotes target network parameters updated via Polyak
averaging: $\bar{\theta} \leftarrow \tau\,\theta + (1-\tau)\,\bar{\theta}$
with $\tau \ll 1$.  The minimum over twin critics addresses the overestimation
bias inherent in Q-learning~\cite{fujimoto2018td3}.

\paragraph{Actor network.}
The policy $\pi_\phi$ is parameterised as a squashed Gaussian: a neural
network outputs mean $\mu_\phi(s)$ and log-standard-deviation
$\log\sigma_\phi(s)$, from which actions are sampled via the reparameterisation
trick:
\begin{equation}
  a = \tanh\!\bigl(\mu_\phi(s) + \sigma_\phi(s) \odot \epsilon\bigr),
  \quad \epsilon \sim \mathcal{N}(0, I)
  \label{eq:sac-reparameterisation}
\end{equation}
The $\tanh$ squashing bounds actions to $[-1, 1]$, which are then rescaled
to the environment's action range $[0, 1]$.  The actor is updated by
minimising:
\begin{equation}
  \mathcal{L}_\pi(\phi) = \mathbb{E}_{s \sim \mathcal{D}}\!\left[
    \alpha \log \pi_\phi(a|s)
    - \min_{j=1,2} Q_{\theta_j}(s, a)
  \right], \quad a \sim \pi_\phi(\cdot|s)
  \label{eq:sac-actor-loss}
\end{equation}

SAC is well-suited to building energy control for several reasons:
(i)~off-policy learning enables sample-efficient training from a replay buffer,
critical when each episode represents a full year of building operation;
(ii)~the continuous stochastic policy naturally handles the continuous cooling
action space; and
(iii)~entropy regularisation prevents the policy from collapsing to a single
cooling strategy, enabling discovery of time-varying control patterns.

% -----------------------------------------------------------------------------
\subsection{CSAC-LB (Constrained SAC with Log-Barrier)}
\label{sec:temp-csac}

The CSAC-LB algorithm, originally proposed by Zhang
et~al.~\cite{zhang2024csaclb} for safe heat pump control, extends SAC with a
constraint-handling mechanism based on interior-point barrier functions from
convex optimisation.  We adopt this framework and adapt it to the CityLearn
cooling case study.

\subsubsection{The Constrained RL Problem}

Standard SAC (Section~\ref{sec:temp-sac-background}) maximises the reward
objective without any notion of safety constraints.  However, the CMDP
formulation (Section~\ref{sec:temp-formulation}) requires that the cumulative
cost remain below a threshold~$d$.  The constrained optimisation problem is:
\begin{equation}
  \max_\pi\; J_R(\pi) \quad \text{subject to} \quad J_C(\pi) \leq d
  \label{eq:cmdp-problem}
\end{equation}
where $J_R(\pi) = \mathbb{E}_\pi\!\left[\sum_t \gamma^t R_t\right]$ is the
discounted return and
$J_C(\pi) = \mathbb{E}_\pi\!\left[\sum_t \gamma^t C_t\right]$ is the
discounted cumulative cost.  The challenge is to enforce this constraint
during online training while maintaining stable policy improvement.

\subsubsection{SAC-Lag: The Lagrangian Approach and Its Limitations}

The standard approach to constrained SAC is SAC-Lag~\cite{ha2020learning,
achiam2021exploration}, which converts the constrained problem into an
unconstrained min-max game using a Lagrange multiplier $\beta \geq 0$:
\begin{equation}
  \max_\pi \min_{\beta \geq 0}\; L(\pi, \beta) \doteq
    f(\pi) - \beta\, g(\pi)
  \label{eq:sac-lag-objective}
\end{equation}
where the reward objective and cost constraint function are:
\begin{align}
  f(\pi) &= \sum_{t=0}^{T} \mathbb{E}_{a_t \sim \pi(\cdot|s_t)}
    \!\left[\gamma^t R(s_t, a_t) + \alpha\, H\!\left(\pi(\cdot|s_t)\right)\right]
  \label{eq:reward-objective} \\
  g(\pi) &= \sum_{t=0}^{T} \mathbb{E}_{a_t \sim \pi(\cdot|s_t)}
    \!\left[\gamma^t \bigl(C(a_t, s_t) - d\bigr)\right]
  \label{eq:cost-constraint}
\end{align}
Here $d$ is the per-step cost limit and $H(\cdot)$ is the entropy term with
temperature $\alpha$.

SAC-Lag updates $\beta$ via dual gradient descent:
\begin{equation}
  J(\beta) = \mathbb{E}_{s_t \sim \mathcal{D}, a_t \sim \pi_\phi(s_t)}
    \!\left[\beta\bigl(d - Q_{\theta_c}(s_t, a_t)\bigr)\right]
  \label{eq:lagrange-update}
\end{equation}
and the actor loss becomes:
\begin{equation}
  J(\phi) = \mathbb{E}_{s_t \sim \mathcal{D}, a_t \sim \pi_\phi(s_t)}
    \!\left[
      \alpha \log \pi_\phi(a_t | s_t)
      - Q_{\theta_r}(s_t, a_t)
      + \beta\, Q_{\theta_c}(s_t, a_t)
    \right]
  \label{eq:sac-lag-actor}
\end{equation}

\paragraph{Failure modes in building control.}
While SAC-Lag is theoretically sound, Zhang et~al.~\cite{zhang2024csaclb}
identify two critical failure modes that make it unreliable for building
energy management:

\begin{enumerate}
  \item \textbf{Multiplier oscillation.}  The dual gradient descent dynamics
    create an adversarial game between the actor (which tries to reduce cost
    to lower $\beta$'s influence) and the multiplier (which grows when cost
    exceeds~$d$).  In practice, this produces oscillatory behaviour:
    $\beta$ overshoots $\to$ policy becomes overly conservative $\to$ cost
    drops well below $d$ $\to$ $\beta$ decreases $\to$ policy becomes
    aggressive $\to$ cost exceeds $d$ again.  This cycle repeats throughout
    training, preventing convergence to a stable operating point near the
    constraint boundary.

  \item \textbf{Data inefficiency from overshooting.}  When $\beta$ overshoots,
    the policy is pushed far into the feasible interior, wasting exploration
    budget on states that are overly conservative.  In building control, where
    each training episode represents a full year of operation and data
    collection is expensive, this inefficiency is particularly costly.  The
    policy spends significant training time learning to return from overly
    conservative regions rather than exploring near the optimal constraint
    boundary.

  \item \textbf{Sensitivity to learning rates.}  The stability of dual gradient
    descent depends critically on the relative learning rates of the actor and
    multiplier.  If the multiplier learning rate is too high, oscillation
    worsens; if too low, constraint violations persist for many episodes before
    correction.  This introduces an additional hyperparameter that requires
    careful tuning for each building control scenario.
\end{enumerate}

These failure modes motivate the replacement of the Lagrangian multiplier with
a barrier function that embeds the constraint directly in the actor loss,
eliminating the dual variable and its associated instabilities.

\subsubsection{Linear Smoothed Log Barrier Function}

CSAC-LB replaces the Lagrangian multiplier with a differentiable barrier
function embedded directly in the actor loss.  This approach draws on
\emph{interior-point methods} from convex optimisation, where log-barrier
functions are used to enforce inequality constraints without introducing
dual variables.

\paragraph{Classical log barrier.}
The standard log barrier for the constraint $x \leq 0$ is
$-\frac{1}{\mu}\log(-x)$, which approaches infinity as $x \to 0^-$ and
is undefined for $x > 0$.  While this prevents the optimiser from entering
the infeasible region, the undefined gradient in the infeasible region
($x > 0$) is problematic for RL: if the policy ever violates the constraint
(which is common during early training), the barrier provides no gradient
signal to guide the policy back to feasibility.

\paragraph{Smoothed extension.}
The \emph{linear smoothed} variant~\cite{zhang2024csaclb} extends the barrier
to the full real line by replacing the logarithmic branch with a linear
extension in the infeasible region:
\begin{equation}
  \bar{\psi}(x) =
  \begin{cases}
    -\dfrac{1}{\mu} \log(-x) & \text{if } x \leq -\dfrac{1}{\mu^2}
    \\[8pt]
    \mu\, x - \dfrac{1}{\mu}\log\!\left(\dfrac{1}{\mu^2}\right)
      + \dfrac{1}{\mu}
    & \text{otherwise}
  \end{cases}
  \label{eq:smoothed-barrier}
\end{equation}
where $\mu > 0$ controls the barrier steepness.

\paragraph{Theoretical properties.}
The smoothed barrier $\bar{\psi}(x)$ has three critical properties that make it
suitable for constrained RL~\cite{zhang2024csaclb}:

\begin{enumerate}
  \item \textbf{Continuity and differentiability.}  The two branches meet at
    the transition point $x = -1/\mu^2$ with matching function values and
    first derivatives.  This ensures smooth gradient flow through the barrier
    during backpropagation, avoiding the gradient discontinuities that plague
    clipped penalty methods.

  \item \textbf{Non-zero gradient everywhere.}  In the feasible region
    ($x \ll 0$), the gradient is $\bar{\psi}'(x) = -1/(\mu x)$, which
    increases as the constraint boundary is approached.  In the infeasible
    region ($x > -1/\mu^2$), the gradient is the constant $\mu$, providing a
    persistent ``recovery signal'' that guides the policy back toward
    feasibility.  This contrasts with clipped barriers where the gradient
    vanishes in the infeasible region, leaving the policy with no directional
    information for constraint satisfaction.

  \item \textbf{Controllable steepness.}  The parameter $\mu$ controls the
    sharpness of the transition between the feasible and infeasible regions.
    As $\mu \to \infty$, $\bar{\psi}(x)$ converges to the indicator function
    $\mathbf{1}[x > 0] \cdot \infty$, recovering the hard constraint.  Smaller
    $\mu$ values create a smoother transition that is easier to optimise but
    allows larger constraint violations.  This provides a principled
    \emph{constraint strictness knob} for the practitioner.
\end{enumerate}

Figure~\ref{fig:barrier-function} illustrates the barrier shape for different
$\mu$ values, showing the smooth transition from logarithmic to linear
behaviour.

\begin{figure}[htbp]
  \centering
  \includegraphics[width=0.45\textwidth]{figures/fig_barrier_function.pdf}
  \caption{Linear smoothed log barrier function $\bar\psi(x)$ for different
    steepness parameters $\mu$.  The dashed line is the indicator function.
    Larger $\mu$ approximates the hard constraint more closely while
    maintaining differentiability everywhere.  The transition point at
    $x = -1/\mu^2$ (where the logarithmic branch meets the linear extension)
    moves closer to zero as $\mu$ increases.  Adapted from
    Zhang et~al.~\cite{zhang2024csaclb}.}
  \label{fig:barrier-function}
\end{figure}

\paragraph{Incorporating the cost limit.}
To apply the barrier to the CMDP cost constraint $Q_C(s,a) \leq d$, the
barrier is composed with a ReLU activation and a shift:
\begin{equation}
  \tilde{\psi}^*(x) = \bar{\psi}\!\left(
    \operatorname{ReLU}(x - d) - \frac{1}{\mu^2}
  \right)
  \label{eq:barrier-shifted}
\end{equation}
where $x = Q_C(s,a)$ is the estimated cost Q-value.  The composition has the
following behaviour:
\begin{itemize}
  \item When $Q_C < d$ (constraint satisfied): $\operatorname{ReLU}(Q_C - d) = 0$,
    so the barrier input is $-1/\mu^2 < 0$, yielding a near-zero penalty.
  \item When $Q_C \approx d$ (near the boundary): the penalty begins to
    increase smoothly, creating a ``soft wall'' effect.
  \item When $Q_C \gg d$ (constraint violated): the penalty grows linearly
    with slope~$\mu$, providing a strong gradient signal to reduce the cost.
\end{itemize}
The shift by $-1/\mu^2$ ensures that the barrier activates slightly
\emph{before} the constraint boundary, providing an early warning gradient
that encourages the policy to stay away from the boundary rather than
approaching it exactly.

\subsubsection{CSAC-LB Architecture}

CSAC-LB extends the standard SAC architecture with an additional pair of
Q-networks dedicated to estimating constraint costs.  The complete architecture
comprises six neural networks:

\begin{itemize}
  \item \textbf{Actor} $\pi_\phi$: Squashed Gaussian policy (as in standard SAC).
  \item \textbf{Twin reward critics} $Q_{R,1}(\theta_{r_1})$,
    $Q_{R,2}(\theta_{r_2})$: Estimate the expected discounted return.
  \item \textbf{Twin cost critics} $Q_{C,1}(\theta_{c_1})$,
    $Q_{C,2}(\theta_{c_2})$: Estimate the expected discounted cost.
  \item \textbf{Target networks}: Slow-moving copies of all critic networks,
    updated via Polyak averaging.
\end{itemize}

Each critic network takes the concatenated state-action pair
$[s, a] \in \mathbb{R}^{57}$ as input (54-dimensional observation $+$
3-dimensional action) and outputs a scalar Q-value.  All networks use identical
two-layer MLP architectures with $[256, 256]$ hidden units and ReLU
activations.

\subsubsection{Reward Critic Training}

The reward critic training follows standard SAC~\cite{haarnoja2018sac}.  The
twin reward critics $Q_{R,1}$, $Q_{R,2}$ are trained to minimise the soft
Bellman residual using the \emph{optimistic minimum} for target computation:
\begin{equation}
  y_R = r_t + \gamma\,(1 - d_t)\!\left[
    \min\!\left(Q_{\bar{R},1}(s', a'),\; Q_{\bar{R},2}(s', a')\right)
    - \alpha \log \pi_\phi(a'|s')
  \right]
  \label{eq:reward-target}
\end{equation}
where $a' \sim \pi_\phi(\cdot|s')$ and $\bar{R}$ denotes target network
parameters.  The loss for each reward critic is:
\begin{equation}
  \mathcal{L}_{Q_R}(\theta_{r_j}) = \mathbb{E}_{(s,a,r,s') \sim \mathcal{D}}
    \!\left[\bigl(Q_{R,j}(s,a) - y_R\bigr)^2\right]
  \label{eq:reward-critic-loss}
\end{equation}
The $\min$ in~\eqref{eq:reward-target} addresses the overestimation bias of
Q-learning: by taking the lower of two estimates, the target is more
conservative, leading to more stable training.

\subsubsection{Cost Critic Training: Pessimistic Estimation}

The cost critic training is the key architectural difference from standard
SAC.  While the reward critic uses $\min$ (optimistic for reward maximisation),
the cost critic uses $\max$ (\emph{pessimistic} for safety):
\begin{equation}
  Q_C(s,a) = \max\!\left(Q_{C,1}(s,a),\; Q_{C,2}(s,a)\right)
  \label{eq:cost-critic-max}
\end{equation}

The cost critic target uses the same pessimistic maximum for bootstrapping:
\begin{equation}
  y_C = c_t + \gamma\,(1 - d_t)\,\max\!\left(
    Q_{\bar{C},1}(s', a'),\; Q_{\bar{C},2}(s', a')
  \right)
  \label{eq:cost-target}
\end{equation}
Note that the cost target does \emph{not} include the entropy term
$-\alpha \log \pi$: entropy regularisation encourages exploration, which is
desirable for reward maximisation but not for cost estimation, where accurate
prediction of constraint violation is the priority.

The loss for each cost critic is:
\begin{equation}
  \mathcal{L}_{Q_C}(\theta_{c_j}) = \mathbb{E}_{(s,a,c,s') \sim \mathcal{D}}
    \!\left[\bigl(Q_{C,j}(s,a) - y_C\bigr)^2\right]
  \label{eq:cost-critic-loss}
\end{equation}

\paragraph{Why pessimistic cost estimation?}
The choice of $\max$ over $\min$ for the cost critic is a critical design
decision motivated by the asymmetry of safety constraints.  In building
control, \emph{underestimating} the cost is dangerous: if the cost critic
predicts that a cooling action is safe when it actually leads to a comfort
violation, the barrier penalty will be too low and the policy will not receive
sufficient gradient to avoid the violation.  By taking the maximum of two
cost estimates, the pessimistic critic errs on the side of caution---it may
occasionally overestimate costs (leading to slightly conservative cooling),
but it will rarely underestimate them (which would lead to undetected
violations).  Zhang et~al.~\cite{zhang2024csaclb} note that this is
particularly important because most transitions in the replay buffer are
constraint-satisfying, creating a systematic underestimation bias that the
pessimistic estimator counteracts.

\subsubsection{Actor Loss with Barrier Penalty}

The actor loss combines the standard SAC reward objective with the barrier
penalty, replacing the Lagrange multiplier entirely:
\begin{equation}
  \mathcal{L}_\pi(\phi) = \mathbb{E}_{s \sim \mathcal{D}}\!\left[
    \underbrace{\alpha \log \pi_\phi(a|s)
      - \min(Q_{R,1}, Q_{R,2})}_{\text{SAC reward maximisation}}
    + \underbrace{\tilde{\psi}^*\!\left(
      \max(Q_{C,1}, Q_{C,2})\right)}_{\text{barrier safety penalty}}
  \right]
  \label{eq:csac-loss}
\end{equation}
where $a \sim \pi_\phi(\cdot|s)$ is sampled via the reparameterisation trick.

The asymmetry between the reward and cost terms is a key design choice:
\begin{itemize}
  \item \textbf{Reward:} $\min(Q_{R,1}, Q_{R,2})$ (optimistic for the actor,
    conservative estimate of reward $\to$ avoids overestimation).
  \item \textbf{Cost:} $\max(Q_{C,1}, Q_{C,2})$ (pessimistic for the actor,
    conservative estimate of cost $\to$ avoids underestimation of danger).
\end{itemize}
This dual asymmetry ensures that the policy is simultaneously optimistic about
rewards (encouraging exploration of high-reward regions) and pessimistic about
costs (discouraging exploration of high-cost regions).

\paragraph{Gradient flow.}
The barrier penalty gradient flows through the cost critic into the actor:
\begin{equation}
  \frac{\partial \mathcal{L}_\pi}{\partial \phi} \ni
    \bar{\psi}'\!\bigl(\text{slack}\bigr) \cdot
    \frac{\partial\, \max(Q_{C,1}, Q_{C,2})}{\partial a} \cdot
    \frac{\partial a}{\partial \phi}
  \label{eq:barrier-gradient}
\end{equation}
When the cost Q-value is well below the limit ($Q_C \ll d$), the barrier
gradient $\bar{\psi}'$ is near zero, and the actor is free to optimise the
reward.  As $Q_C \to d$, the barrier gradient increases, progressively
steering the actor away from constraint-violating actions.  This creates
a smooth transition from reward-driven to safety-driven behaviour, unlike
SAC-Lag where the Lagrange multiplier applies a constant scaling regardless of
proximity to the constraint boundary.

\subsubsection{Training Procedure}

Algorithm~\ref{alg:csac-lb} presents the complete CSAC-LB training loop.
The key differences from standard SAC are highlighted: the cost critic update
with pessimistic targets (line~\ref{line:cost-target}), and the barrier penalty
in the actor loss (line~\ref{line:barrier-penalty}).

\begin{algorithm}[htbp]
\caption{CSAC-LB Training (adapted from Zhang et~al.~\cite{zhang2024csaclb})}
\label{alg:csac-lb}
\begin{algorithmic}[1]
\REQUIRE Reward critics $Q_{R,1}, Q_{R,2}$; cost critics $Q_{C,1}, Q_{C,2}$;
  actor $\pi_\phi$; target networks $\bar{Q}_{R}, \bar{Q}_{C}$;
  replay buffer $\mathcal{D}$; cost limit $d$; barrier factor $\mu$
\STATE Initialise all networks with Kaiming uniform weights
\STATE $\text{barrier\_shift} \gets 1/\mu^2$
\FOR{each environment step $t = 1, 2, \ldots, T_{\max}$}
  \IF{$t \leq t_{\text{warmup}}$}
    \STATE $a_t \gets \text{UniformRandom}(\mathcal{A})$
  \ELSE
    \STATE $a_t \sim \pi_\phi(\cdot | s_t)$
    \hfill $\triangleright$ \textit{Stochastic action via reparameterisation}
  \ENDIF
  \STATE Execute $a_t$, observe $r_t, c_t, s_{t+1}$
  \STATE Store $(s_t, a_t, r_t, c_t, s_{t+1})$ in $\mathcal{D}$
  \IF{$t > t_{\text{warmup}}$}
    \STATE Sample mini-batch $\mathcal{B} \sim \mathcal{D}$
    \STATE \textbf{// Update reward critics}
    \STATE $a' \sim \pi_\phi(\cdot | s')$
    \STATE $y_R \gets r + \gamma(1-d_t)[\min(\bar{Q}_{R,1}(s',a'), \bar{Q}_{R,2}(s',a')) - \alpha \log \pi_\phi(a'|s')]$
    \STATE Update $Q_{R,1}, Q_{R,2}$ to minimise $\|Q_{R,j}(s,a) - y_R\|^2$
    \STATE \textbf{// Update cost critics (pessimistic)} \label{line:cost-target}
    \STATE $y_C \gets c + \gamma(1-d_t)\,\max(\bar{Q}_{C,1}(s',a'), \bar{Q}_{C,2}(s',a'))$
    \hfill $\triangleright$ \textit{max, not min}
    \STATE Update $Q_{C,1}, Q_{C,2}$ to minimise $\|Q_{C,j}(s,a) - y_C\|^2$
    \IF{$t \bmod \text{policy\_delay} = 0$}
      \STATE \textbf{// Update actor with barrier penalty} \label{line:barrier-penalty}
      \STATE $\tilde{a} \sim \pi_\phi(\cdot | s)$
      \STATE $Q_R \gets \min(Q_{R,1}(s,\tilde{a}),\; Q_{R,2}(s,\tilde{a}))$
      \STATE $Q_C \gets \max(Q_{C,1}(s,\tilde{a}),\; Q_{C,2}(s,\tilde{a}))$
      \STATE $\text{slack} \gets \operatorname{ReLU}(Q_C - d) - \text{barrier\_shift}$
      \STATE $\mathcal{L}_\pi \gets \alpha \log \pi_\phi(\tilde{a}|s) - Q_R + \bar{\psi}(\text{slack})$
      \STATE Update $\phi$ to minimise $\mathcal{L}_\pi$
      \STATE \textbf{// Polyak-average all target networks}
      \STATE $\bar{\theta} \gets \tau\,\theta + (1-\tau)\,\bar{\theta}$
      \hfill $\triangleright$ \textit{For reward, cost critics and actor}
    \ENDIF
  \ENDIF
\ENDFOR
\end{algorithmic}
\end{algorithm}

\paragraph{Key design choices.}
Several aspects of the training procedure merit discussion:

\begin{itemize}
  \item \textbf{Policy delay} (update actor every 2 critic steps): This gives
    the critic networks time to converge to more accurate Q-estimates before
    the actor uses them, reducing variance in the policy gradient.  This is
    particularly important for CSAC-LB because the actor loss depends on
    \emph{both} the reward and cost critics---inaccurate cost estimates early
    in training could cause the barrier to fire inappropriately.

  \item \textbf{Fixed entropy temperature} ($\alpha = 0.2$): Unlike many SAC
    implementations that use automatic entropy tuning (adjusting $\alpha$ to
    maintain a target entropy $-\dim(\mathcal{A})$), we fix $\alpha = 0.2$
    following Zhang et~al.~\cite{zhang2024csaclb}.  In the constrained
    setting, automatic entropy tuning can interfere with the barrier mechanism:
    when the barrier penalty increases (pushing the actor toward conservative
    actions), the entropy may decrease, triggering $\alpha$ to increase, which
    counteracts the barrier's effect.

  \item \textbf{Uniform replay buffer}: We use standard uniform sampling rather
    than prioritised experience replay.  While prioritised replay can
    accelerate learning in unconstrained settings, it can bias the cost critic
    by over-representing rare constraint-violating transitions, potentially
    leading to overestimation of costs.

  \item \textbf{Gradient clipping} (max norm $= 40.0$): Applied to all
    networks to prevent large barrier gradients from destabilising training
    when the cost critic predicts values far above the limit during early
    epochs.
\end{itemize}

\subsubsection{Architecture and Hyperparameters}

Table~\ref{tab:csac-hyperparams} summarises the CSAC-LB hyperparameters.
The network architecture and learning rates follow the configuration of
Zhang et~al.~\cite{zhang2024csaclb}; the cost limit, barrier factor, and
reward weights are adapted for the CityLearn episode structure and comfort
requirements.

\begin{table}[htbp]
\centering
\caption{CSAC-LB hyperparameters for the temperature case study.}
\label{tab:csac-hyperparams}
\begin{tabular}{lcc}
\toprule
\textbf{Hyperparameter} & \textbf{Value} & \textbf{Source} \\
\midrule
\multicolumn{3}{l}{\textit{Network architecture}} \\
Hidden layers (actor \& all critics) & $[256, 256]$ & \cite{zhang2024csaclb} \\
Activation function & ReLU & \cite{zhang2024csaclb} \\
Weight initialisation & Kaiming uniform & \cite{zhang2024csaclb} \\
Output squashing (actor) & $\tanh$ & \cite{haarnoja2018sac} \\
\midrule
\multicolumn{3}{l}{\textit{Optimisation}} \\
Learning rate (all networks) & $10^{-3}$ & \cite{zhang2024csaclb} \\
Optimiser & Adam & \cite{zhang2024csaclb} \\
Batch size & 256 & \cite{zhang2024csaclb} \\
Discount factor $\gamma$ & 0.99 & \cite{zhang2024csaclb} \\
Polyak averaging $\tau$ & 0.005 & \cite{zhang2024csaclb} \\
Policy delay & 2 & \cite{zhang2024csaclb} \\
Gradient clipping (max norm) & 40.0 & \cite{zhang2024csaclb} \\
\midrule
\multicolumn{3}{l}{\textit{Replay buffer}} \\
Buffer capacity & $3 \times 10^6$ & Tuned \\
Random exploration steps & 100 & Tuned \\
Update frequency & 1 (every env step) & \cite{zhang2024csaclb} \\
\midrule
\multicolumn{3}{l}{\textit{Constraint handling}} \\
Cost limit $d$ & 0.5 & Tuned \\
Barrier steepness $\mu$ & 25.0 & Tuned \\
Barrier shift $1/\mu^2$ & 0.0016 & Derived \\
\midrule
\multicolumn{3}{l}{\textit{Entropy and reward}} \\
Entropy temperature $\alpha$ & 0.2 (fixed) & \cite{zhang2024csaclb} \\
Comfort reward weight $\lambda_{\text{indoor}}$ & 1.2 & Tuned \\
\midrule
\multicolumn{3}{l}{\textit{Training schedule}} \\
Steps per epoch & 2,208 & Environment \\
Total training epochs & 40 & Tuned \\
Total training steps & 88,320 & Derived \\
\bottomrule
\end{tabular}
\end{table}

\paragraph{Hyperparameter rationale.}
The cost limit $d = 0.5$ and barrier steepness $\mu = 25$ were selected via
grid search over $d \in \{0.1, 0.5, 1.0, 2.0, 5.0\}$ and
$\mu \in \{5, 10, 15, 25, 50\}$.  Low $d$ values ($d = 0.1$) produced overly
conservative policies that sacrificed energy efficiency for minimal violation
rates, while high values ($d = 5.0$) resulted in insufficient constraint
enforcement.  The combination $d = 0.5$, $\mu = 25$ achieved the best
trade-off: the policy reaches violation rates below $11\%$ while consuming
less energy than the bang-bang RBC baseline.

The comfort reward weight $\lambda_{\text{indoor}} = 1.2$ serves a dual
purpose: it provides a shaped reward gradient for comfort improvement
(complementing the binary cost signal $C_t$), and it scales the comfort
component to be comparable in magnitude to the other STEMS reward components.
Without this weight, the comfort reward would be dominated by the economic and
stability terms, and the policy would learn to ignore temperature deviations
until the barrier penalty fires---at which point the correction is abrupt
rather than gradual.

% -----------------------------------------------------------------------------
\subsection{Baseline: Rule-Based Cooling Controller (RBC)}
\label{sec:temp-rbc}

The primary baseline is a deterministic \emph{bang-bang} thermostat controller,
representative of the on/off control logic found in commodity and legacy HVAC
systems.  The controller applies a simple binary decision rule:
\begin{equation}
  a_t^{(i)} =
  \begin{cases}
    1.0 & \text{if } T_t^{(i)} > T_{\text{bb}} \\
    0.0 & \text{otherwise}
  \end{cases}
  \label{eq:bangbang-rbc}
\end{equation}
where $T_{\text{bb}} = 24\,^\circ$C is the switching threshold, set
$2\,^\circ$C below the upper comfort bound of $26\,^\circ$C.

This controller exhibits three characteristic inefficiencies that are
well-documented in building energy literature:
\begin{enumerate}
  \item \textbf{Energy waste from premature activation.}  The threshold at
    $24\,^\circ$C triggers full cooling for buildings that are still within the
    comfort band, consuming energy unnecessarily.
  \item \textbf{Grid instability from hard switching.}  The binary on/off
    transitions create large load spikes every time a building crosses the
    threshold, resulting in high ramping rates ($2.322$ KPI, the worst among
    all baselines).
  \item \textbf{Poor comfort despite high energy use.}  The all-or-nothing
    response causes temperature oscillation around the threshold: full cooling
    drives the temperature below $24\,^\circ$C, the controller switches off,
    the building heats back up, and the cycle repeats---yielding a $51.6\%$
    violation rate despite a mean cooling intensity of $0.207$.
\end{enumerate}

A \emph{Zero} baseline (no cooling action) establishes the worst-case comfort
floor.


% =============================================================================
\section{Evaluation Metrics}
\label{sec:temp-metrics}

Following the evaluation methodology of STEMS~\cite{zhang2025stems}, we
evaluate all methods using the following key performance metrics.

\paragraph{Comfort violation rate.}
The fraction of active time steps at which any building's indoor temperature
falls outside the comfort band:
\begin{equation}
  V_{\text{rate}} = \frac{1}{T_{\text{active}}}
    \sum_{t=1}^{T_{\text{active}}} \mathbf{1}\!\left[
      \exists\, i : T_t^{(i)} < 20\,^\circ\text{C}
      \;\text{or}\; T_t^{(i)} > 26\,^\circ\text{C}
    \right]
  \label{eq:violation-rate}
\end{equation}

\paragraph{Discomfort proportion.}
The CityLearn KPI measuring the fraction of building-hour pairs outside the
comfort band, accounting for both hot and cold discomfort:
\begin{equation}
  D_{\text{prop}} = \frac{1}{T \cdot N}
    \sum_{t=1}^{T} \sum_{i=1}^{N}
    \mathbf{1}\!\left[T_{t}^{(i)} < 20\,^\circ\text{C}
      \;\text{or}\; T_{t}^{(i)} > 26\,^\circ\text{C}\right]
  \label{eq:discomfort-proportion}
\end{equation}

\paragraph{Occupancy-weighted discomfort rate.}
The fraction of occupied building-hours that experience discomfort:
\begin{equation}
  D_{\text{occ}} = \frac{\text{discomfort\_count}}{\text{occupied\_count}}
  \label{eq:discomfort-occ-rate}
\end{equation}

\paragraph{Cost and emissions.}
Total electricity cost and carbon emissions, reported as CityLearn KPI ratios
normalised relative to a no-storage baseline following~\cite{zhang2025stems}.
Values below 1.0 indicate improvement over the no-storage baseline.

\paragraph{Grid stability metrics.}
\begin{itemize}
  \item \textbf{Ramping rate:} average absolute change in grid load between
        consecutive time steps, normalised to the RBC baseline.
  \item \textbf{Daily peak:} average maximum daily power demand, normalised.
  \item \textbf{Load factor:} daily one-minus-load-factor, measuring demand
        variability.
\end{itemize}

\paragraph{Mean cooling action.}
Average cooling intensity across all buildings and time steps, indicating
energy aggressiveness of the policy.


% =============================================================================
\section{Results and Discussion}
\label{sec:temp-results}

% -----------------------------------------------------------------------------
\subsection{Training Dynamics}
\label{sec:temp-results-training}

Figure~\ref{fig:training-curves} shows the CSAC-LB training trajectory over
40 epochs.  The agent rapidly reduces the episode comfort cost from $\sim$290
to $\sim$1.4 within the first 10 epochs, while the episode return initially
drops (due to aggressive cooling) but recovers as the policy learns more
efficient control.  The barrier penalty and cost critic Q-value converge in
tandem, confirming that the barrier mechanism guides exploration toward the
feasible boundary.

\begin{figure}[htbp]
  \centering
  \includegraphics[width=\textwidth]{figures/fig_training_curves.png}
  \caption{CSAC-LB training dynamics over 40 epochs.
    (a)~Episode return: starts lower due to aggressive cooling but improves
    as the policy learns efficiency.
    (b)~Comfort violation cost: rapidly driven near zero.
    (c)~Actor loss convergence.
    (d)~Barrier penalty and cost Q-value, showing the barrier mechanism
    converging.}
  \label{fig:training-curves}
\end{figure}

Figure~\ref{fig:training-energy-comfort} compares the training trajectories
against the bang-bang RBC baseline (dashed black), following the visualisation
style of Zhang et~al.~\cite{zhang2024csaclb}.  CSAC-LB achieves a dramatic
cost reduction early in training, consistent with the original CSAC-LB
findings that the log-barrier method guides exploration toward the constraint
boundary more effectively than Lagrangian approaches.

\begin{figure}[htbp]
  \centering
  \includegraphics[width=\textwidth]{figures/fig_training_energy_comfort.png}
  \caption{Episode return and comfort violation cost during training, with
    bang-bang RBC baseline (dashed).
    (a)~CSAC-LB trades lower initial return for constraint satisfaction;
    return recovers as the policy learns efficient cooling.
    (b)~CSAC-LB reduces comfort violations by an order of magnitude below
    the RBC baseline.}
  \label{fig:training-energy-comfort}
\end{figure}

Figure~\ref{fig:pareto-front} visualises the reward--cost Pareto front during
training, inspired by the exploration analysis of
Zhang et~al.~\cite{zhang2024csaclb} (their Fig.~5).  Each point represents one
epoch's evaluation, coloured by training progress.  CSAC-LB explores along the
constraint boundary (low cost, improving return), moving away from the
bang-bang RBC operating point toward a superior reward--cost trade-off.

\begin{figure}[htbp]
  \centering
  \includegraphics[width=0.75\textwidth]{figures/fig_pareto_front.png}
  \caption{Pareto front of episode return vs comfort violation cost.
    CSAC-LB (circles, coloured by epoch) explores along the low-cost
    boundary.  The bang-bang RBC (black cross) occupies a high-cost,
    low-return region.  Stars mark the best and final CSAC-LB checkpoints.}
  \label{fig:pareto-front}
\end{figure}

% -----------------------------------------------------------------------------
\subsection{Long-Term Performance Comparison}
\label{sec:temp-results-longterm}

Table~\ref{tab:temp-main-results} presents the full-year evaluation results
for all methods.  All metrics are computed from a single deterministic rollout
of each converged policy over the complete $2{,}207$-step episode.
Figure~\ref{fig:eval-bar-chart} provides a visual comparison across six key
metrics.

\begin{figure}[htbp]
  \centering
  \includegraphics[width=\textwidth]{figures/fig_eval_bar_chart.png}
  \caption{Full-year evaluation comparison.  CSAC-LB (red) achieves
    dramatically lower violation rates while also outperforming the RBC
    baseline on all energy efficiency metrics.}
  \label{fig:eval-bar-chart}
\end{figure}

\begin{table}[htbp]
\centering
\caption{Temperature control case study: full-year evaluation results.
  Violation rate and discomfort proportion are lower-is-better; CityLearn KPI
  ratios (cost, emissions, consumption) are normalised relative to a no-storage
  baseline where values $< 1.0$ indicate improvement.  CSAC-LB results that
  dominate the RBC baseline are in \textbf{bold}.}
\label{tab:temp-main-results}
\begin{tabular}{lcccccc}
\toprule
\textbf{Method}
  & \textbf{Violation}
  & \textbf{Discomfort}
  & \textbf{Cost}
  & \textbf{Emissions}
  & \textbf{Elec.\ Cons.}
  & \textbf{Ramping} \\
  & \textbf{Rate}
  & \textbf{Proportion}
  & \textbf{(KPI)}
  & \textbf{(KPI)}
  & \textbf{(KPI)}
  & \textbf{(KPI)} \\
\midrule
Zero (no cooling)
  & 0.995 & 0.962 & 0.527 & 0.562 & 0.559 & 0.905 \\
RBC
  & 0.516 & 0.376 & 1.010 & 1.043 & 1.041 & 2.322 \\
\midrule
CSAC-LB (best)
  & \textbf{0.060} & \textbf{0.308} & 1.012 & 1.061 & 1.058 & \textbf{1.024} \\
CSAC-LB (final)
  & \textbf{0.049} & \textbf{0.211} & \textbf{0.958} & \textbf{0.993} & \textbf{0.991} & \textbf{1.076} \\
\bottomrule
\end{tabular}
\end{table}

\paragraph{CSAC-LB achieves dramatic safety improvement.}
The most striking result is the reduction in comfort violation rate by CSAC-LB.
While the bang-bang RBC violates the comfort band in $51.6\%$ of time steps,
CSAC-LB (final checkpoint) reduces this to $4.9\%$---a $90.5\%$ relative
reduction.  The best-epoch checkpoint achieves $6.0\%$, representing an
$88.4\%$ improvement.  This demonstrates that the smoothed log-barrier penalty
effectively enforces the comfort constraint during training, producing a policy
that maintains indoor temperatures within $[20\,^\circ\text{C},
26\,^\circ\text{C}]$ for the vast majority of the operating year.

\paragraph{CSAC-LB dominates the RBC baseline across all metrics.}
Unlike typical safety--efficiency trade-offs reported in prior
work~\cite{zhang2024csaclb}, the CSAC-LB final checkpoint achieves
\emph{simultaneous improvement} over the bang-bang RBC on every key metric:
cost ($0.958$ vs $1.010$, $-5.1\%$), emissions ($0.993$ vs $1.043$,
$-4.8\%$), consumption ($0.991$ vs $1.041$, $-4.8\%$), daily peak ($1.071$
vs $1.131$, $-5.3\%$), and ramping ($1.076$ vs $2.322$, $-53.7\%$)---all
while reducing violations by $90.5\%$.  This full dominance is possible
because the bang-bang controller wastes energy through its binary switching:
it applies maximum cooling whenever any building exceeds $24\,^\circ$C, then
switches off completely, causing temperature oscillation and grid load spikes.
CSAC-LB, by contrast, learns \emph{when} and \emph{how much} to cool,
achieving both comfort and efficiency through continuous, modulated control.

\paragraph{Reward decomposition analysis.}
Table~\ref{tab:temp-reward-decomposition} breaks down the cumulative reward
components.  The bang-bang RBC achieves high stability ($1{,}703.6$) and
renewable utilisation ($1{,}136.6$) but incurs a massive comfort penalty
($-3{,}059.7$) due to its oscillating temperature profile.  CSAC-LB (final)
achieves a dramatically better comfort reward ($-71.3$ vs $-3{,}059.7$)
while maintaining comparable economic efficiency ($-246.5$ vs $-229.1$).

\begin{table}[htbp]
\centering
\caption{Cumulative reward component decomposition over the full evaluation
  episode.  Higher (less negative) is better for all components.}
\label{tab:temp-reward-decomposition}
\begin{tabular}{lrrrr}
\toprule
\textbf{Method}
  & $R^{\text{economic}}$
  & $R^{\text{stability}}$
  & $R^{\text{renewable}}$
  & $R^{\text{comfort}}$ \\
\midrule
RBC
  & $-229.1$ & $1{,}703.6$ & $1{,}136.6$ & $-3{,}059.7$ \\
CSAC-LB (best)
  & $-264.1$ & $1{,}269.4$ & $905.1$ & $-35.1$ \\
CSAC-LB (final)
  & $-246.5$ & $1{,}619.2$ & $956.7$ & $-71.3$ \\
\bottomrule
\end{tabular}
\end{table}


% -----------------------------------------------------------------------------
\subsection{Extreme Weather Stress Test}
\label{sec:temp-results-heatwave}

Following the extreme weather evaluation methodology of
STEMS~\cite{zhang2025stems} (their Section~V-B), we evaluate the robustness
of CSAC-LB under heat-wave conditions.  STEMS evaluates under two extreme
weather scenarios:

\begin{itemize}
  \item \textbf{Heat Wave:} outdoor temperature exceeding $35\,^\circ$C
    (occurring $21.6\%$ of the year in Travis County), with cooling demand
    increased by $150\%$ and solar generation enhanced by $20\%$.
  \item \textbf{Cold Wave:} outdoor temperature below $-3\,^\circ$C
    (occurring $3.0\%$ of the year), with heating demand increased by $180\%$
    and battery capacity reduced by $15\%$.
\end{itemize}

Since our case study uses a cooling-only action interface, we adapt the STEMS
heat-wave protocol while omitting the cold-wave scenario (which targets heating
systems).  Our adaptation applies the following perturbations to identified
heat-wave windows:

\begin{itemize}
  \item Cooling demand multiplied by $2.5\times$ (more aggressive than STEMS's
    $1.5\times$ to stress-test the cooling-only policy),
  \item Solar generation multiplied by $1.2\times$ (matching STEMS),
  \item Windows identified by contiguous periods where outdoor temperature
    exceeds a threshold ($25.7\,^\circ$C for 24h, $25.0\,^\circ$C for 72h).
    The 72h scenario requires a lower threshold because the dataset contains
    no contiguous spell of 72 hours above $25.7\,^\circ$C (the longest such
    spell is only 42\,h); lowering the threshold to $25.0\,^\circ$C yields a
    90-hour spell from which a 72-hour window is extracted.
\end{itemize}

The heat-wave perturbation is applied directly to the CityLearn building energy
simulation arrays: the \texttt{cooling\_demand} and
\texttt{cooling\_demand\_without\_control} arrays are scaled within the window
bounds, and the \texttt{solar\_generation} array is similarly amplified.  Both
CSAC-LB and RBC are evaluated on the \emph{same perturbed environment}
for a fair comparison.

\subsubsection{Heat-Wave Results}

Table~\ref{tab:temp-heatwave} presents results for 24-hour and 72-hour
heat-wave windows.

\begin{table}[htbp]
\centering
\caption{Heat-wave stress test results.  In the 24h scenario, CSAC-LB
  dominates RBC on \emph{all} metrics including cost and emissions.  In
  the 72h scenario, RBC achieves lower energy KPIs but at much higher
  violation rates.  KPI values are normalised relative to a no-storage
  baseline.  Bold indicates the better value per metric.}
\label{tab:temp-heatwave}
\begin{tabular}{llccccc}
\toprule
\textbf{Duration}
  & \textbf{Method}
  & \textbf{Violation}
  & \textbf{Discomfort}
  & \textbf{Cost}
  & \textbf{Emissions}
  & \textbf{Ramping} \\
  & & \textbf{Rate} & \textbf{Prop.} & \textbf{(KPI)} & \textbf{(KPI)} & \textbf{(KPI)} \\
\midrule
\multirow{2}{*}{24\,h}
  & RBC     & 0.700 & 0.250 & 1.089 & 1.088 & 2.060 \\
  & CSAC-LB & \textbf{0.000} & \textbf{0.000} & \textbf{1.034} & \textbf{1.035} & \textbf{1.141} \\
\midrule
\multirow{2}{*}{72\,h}
  & RBC     & 0.517 & 0.385 & \textbf{0.976} & 1.016 & 1.574 \\
  & CSAC-LB & \textbf{0.086} & \textbf{0.185} & 1.022 & \textbf{1.029} & \textbf{1.061} \\
\bottomrule
\end{tabular}
\end{table}

\paragraph{24-hour heat wave.}
CSAC-LB achieves a \emph{perfect} violation rate of $0.0\%$ and zero
discomfort proportion, completely eliminating comfort violations during the
24-hour extreme heat event.  In contrast, the bang-bang RBC suffers a $70.0\%$
violation rate and $25.0\%$ discomfort proportion.  Crucially, CSAC-LB also
achieves lower cost ($1.034$ vs $1.089$, $-5.0\%$), lower emissions ($1.035$
vs $1.088$, $-4.9\%$), and dramatically lower ramping ($1.141$ vs $2.060$,
$-44.6\%$).  This full dominance demonstrates that CSAC-LB's continuous
modulated control outperforms the bang-bang RBC's wasteful binary switching
even under extreme heat conditions.

\paragraph{72-hour heat wave.}
Over the longer 72-hour window, CSAC-LB maintains a violation rate of only
$8.6\%$, compared to RBC's $51.7\%$---an $83.3\%$ relative reduction.  The
ramping rate improves from $1.574$ to $1.061$ ($-32.6\%$).  RBC achieves
a lower cost KPI ($0.976$ vs $1.022$, $4.6\%$ gap), while CSAC-LB has
slightly better emissions ($1.029$ vs $1.016$).  The 72h scenario is more
nuanced than the 24h case: during extended heat, the bang-bang RBC runs at
full power for sustained periods (reducing wasted switching cycles), but still
suffers from cold discomfort ($38.5\%$) due to over-cooling between heat
peaks.  CSAC-LB reduces discomfort to $18.5\%$ and violations by $83\%$, at
a modest energy cost premium.

\subsubsection{Multi-Scenario Radar Chart Analysis}

Figure~\ref{fig:stress-test-radar} presents STEMS-style radar charts comparing
CSAC-LB's improvement over RBC across three scenarios, following the
visualisation approach of Zhang et~al.~\cite{zhang2025stems} (their Fig.~3).

\begin{figure}[htbp]
  \centering
  \includegraphics[width=\textwidth]{figures/fig_stress_test_radar.png}
  \caption{Radar charts showing CSAC-LB's percentage improvement over
    RBC across five technical metrics for full-year, 24h heat-wave, and 72h
    heat-wave scenarios.  Positive values indicate improvement.  CSAC-LB
    achieves positive improvement on all metrics in the full-year and 24h
    scenarios.  In the 72h scenario, cost/emission/consumption metrics are
    negative due to the extended heat period where bang-bang cooling is
    coincidentally energy-efficient.
    Adapted from the STEMS radar chart visualisation~\cite{zhang2025stems}.}
  \label{fig:stress-test-radar}
\end{figure}

The radar charts reveal two distinct patterns.  In the \emph{full-year} and
\emph{24h heat-wave} scenarios, CSAC-LB achieves positive improvement on all
five metrics: ramping reduction of $53.7\%$ (full-year) and $44.6\%$ (24h),
cost reduction of $5.1\%$ (full-year) and $5.0\%$ (24h), and comparable
emission/consumption improvements.  In the \emph{72h heat-wave} scenario,
CSAC-LB shows a modest cost increase ($+4.6\%$) and marginal emission
difference, though ramping improves by $32.6\%$.  The consistent ramping
improvement across all scenarios confirms that CSAC-LB's continuous control
produces smoother grid operation regardless of weather conditions.

\subsubsection{Discomfort and Safety Violation Analysis}

Figure~\ref{fig:stress-test-discomfort} compares discomfort levels and safety
violation rates across scenarios, following the STEMS paper's Figures~4 and~5.

\begin{figure}[htbp]
  \centering
  \includegraphics[width=\textwidth]{figures/fig_stress_test_discomfort_violations.png}
  \caption{Discomfort and safety violation comparison across scenarios,
    following the STEMS evaluation format~\cite{zhang2025stems}.
    (a)~Discomfort proportion: CSAC-LB achieves zero discomfort in the 24h
    heat wave but shows higher discomfort in the 72h case.
    (b)~Safety violations: CSAC-LB reduces violations by $90\%$ (full-year),
    $100\%$ (24h), and $83\%$ (72h) relative to RBC.}
  \label{fig:stress-test-discomfort}
\end{figure}

The safety violation reduction is the most striking result:
\begin{itemize}
  \item \textbf{Full-year:} $51.6\% \to 4.9\%$ ($-90.5\%$)
  \item \textbf{24h heat wave:} $70.0\% \to 0.0\%$ ($-100\%$)
  \item \textbf{72h heat wave:} $51.7\% \to 8.6\%$ ($-83.3\%$)
\end{itemize}
These reductions are comparable to STEMS's reported safety violation reduction
from $35.1\%$ to $5.6\%$ ($-84.0\%$)~\cite{zhang2025stems}, demonstrating
that the log-barrier mechanism achieves CBF-comparable safety enforcement
without requiring the explicit constraint-checking infrastructure of STEMS.

\subsubsection{Hot vs Cold Discomfort Breakdown}

Figure~\ref{fig:hot-cold-breakdown} decomposes the discomfort proportion into
hot ($T > 26\,^\circ$C) and cold ($T < 20\,^\circ$C) components.

\begin{figure}[htbp]
  \centering
  \includegraphics[width=\textwidth]{figures/fig_stress_test_hot_cold_breakdown.png}
  \caption{Hot vs cold discomfort breakdown across scenarios.  The bang-bang
    RBC suffers severe cold discomfort ($0.376$ full-year) from over-cooling,
    while CSAC-LB reduces this to $0.186$.  CSAC-LB introduces modest hot
    discomfort ($0.025$ full-year vs RBC's $0.000$) as it avoids the RBC's
    aggressive full-power cooling.  In the 24h heat-wave, both components
    are zero for CSAC-LB.}
  \label{fig:hot-cold-breakdown}
\end{figure}

This analysis reveals an important distinction between the two controllers'
failure modes.  The bang-bang RBC's dominant failure is \emph{cold} discomfort
($0.376$ full-year): its binary full-power cooling drives temperatures well
below the $20\,^\circ$C lower comfort bound, and the lack of modulation means
the building oscillates between over-cooled and over-heated states.  CSAC-LB
dramatically reduces cold discomfort to $0.186$ ($-50.5\%$) by learning
graded cooling actions, but introduces modest hot discomfort ($0.025$ vs
$0.000$) as it sometimes under-cools to conserve energy.  In the 72h
heat-wave, cold discomfort rises for CSAC-LB ($0.108$) as the sustained heat
event pushes the learned policy toward more aggressive cooling.  Future work
could further reduce the remaining discomfort through:
\begin{itemize}
  \item asymmetric comfort penalties that weight hot and cold violations differently,
  \item separate cost terms for hot and cold violations in the CMDP, or
  \item adaptive barrier annealing that gradually tightens constraints during training.
\end{itemize}

Figure~\ref{fig:stress-test-detailed} provides a comprehensive per-metric
comparison across all three scenarios.

\begin{figure}[htbp]
  \centering
  \includegraphics[width=\textwidth]{figures/fig_stress_test_detailed_metrics.png}
  \caption{Detailed metric comparison across full-year, 24h, and 72h heat-wave
    scenarios.  Top row: violation rate, cost, emissions.  Bottom row:
    electricity consumption, ramping, discomfort rate.}
  \label{fig:stress-test-detailed}
\end{figure}


% -----------------------------------------------------------------------------
\subsection{CityLearn KPI Comparison}
\label{sec:temp-results-kpis}

Table~\ref{tab:temp-kpis} presents a broader set of CityLearn KPIs comparing
RBC and CSAC-LB (final checkpoint).

\begin{table}[htbp]
\centering
\caption{Extended CityLearn KPI comparison.  All values are normalised ratios
  where $<1.0$ is better (except thermal resilience, where lower
  one-minus-resilience is better).}
\label{tab:temp-kpis}
\begin{tabular}{lcc}
\toprule
\textbf{KPI}
  & \textbf{RBC}
  & \textbf{CSAC-LB (final)} \\
\midrule
Cost total                & 1.010 & \textbf{0.958} \\
Carbon emissions          & 1.043 & \textbf{0.993} \\
Electricity consumption   & 1.041 & \textbf{0.991} \\
Daily peak                & 1.131 & \textbf{1.071} \\
Ramping average           & 2.322 & \textbf{1.076} \\
Discomfort proportion     & 0.376 & \textbf{0.211} \\
Discomfort (hot)          & \textbf{0.000} & 0.025 \\
Discomfort (cold)         & 0.376 & \textbf{0.186} \\
Daily load factor (1$-$LF)& 1.101 & \textbf{0.998} \\
Zero net energy           & 1.008 & \textbf{0.984} \\
\bottomrule
\end{tabular}
\end{table}

CSAC-LB (final) achieves lower values than the bang-bang RBC on 9 out of 10
KPIs, demonstrating comprehensive improvement across cost, emissions,
consumption, peak demand, ramping, load factor, and net energy metrics.  The
only KPI where the bang-bang RBC performs better is hot discomfort proportion
($0.000$ vs $0.025$), because its aggressive full-power cooling during hot
periods eliminates hot violations entirely---at the cost of extreme cold
discomfort ($0.376$) from over-cooling and subsequent temperature oscillation.


% -----------------------------------------------------------------------------
\subsection{Discussion}
\label{sec:temp-discussion}

Figure~\ref{fig:radar-chart} presents a STEMS-style radar chart showing
the percentage improvement of CSAC-LB over the bang-bang RBC across five
key metrics, following the visualisation approach of
Zhang et~al.~\cite{zhang2025stems} (their Fig.~3).

\begin{figure}[htbp]
  \centering
  \includegraphics[width=0.65\textwidth]{figures/fig_radar_chart.png}
  \caption{Radar chart showing CSAC-LB percentage improvement over the
    bang-bang RBC.  Positive values indicate better performance.  The
    CSAC-LB final checkpoint achieves positive improvement on all five
    axes, confirming full dominance over the baseline.}
  \label{fig:radar-chart}
\end{figure}

\subsubsection{Comparison with Original CSAC-LB Findings}

Our results are consistent with the findings of Zhang
et~al.~\cite{zhang2024csaclb} on the I4B heat pump simulator:

\begin{itemize}
  \item \textbf{Boundary exploration:} The Pareto front
    (Figure~\ref{fig:pareto-front}) confirms CSAC-LB's characteristic
    behaviour of exploring along the feasible boundary, matching the original
    paper's Figure~5 where CSAC-LB maintains closer proximity to the optimal
    constraint boundary compared to SAC-Lag and CPO.

  \item \textbf{Comfort--energy trade-off:} Zhang et~al.\ report that CSAC-LB
    achieves low average temperature deviation (0.047\,K for Building~1) with
    moderate energy use (3,652\,kWh), outperforming SAC-Lag which shows higher
    energy consumption (3,584\,kWh) with worse deviation (0.199\,K).  Our
    results extend this finding: CSAC-LB achieves far lower violation rates
    than the bang-bang RBC \emph{and} lower energy consumption, demonstrating
    that the learned policy achieves both safety and efficiency.

  \item \textbf{Training stability:} Our training curves
    (Figure~\ref{fig:training-curves}) show stable CSAC-LB convergence,
    matching the original paper's observation that CSAC-LB avoids the training
    instabilities exhibited by SAC-Lag and CPO.
\end{itemize}

\subsubsection{Key Insights}

The temperature control case study reveals three key insights:

\begin{enumerate}
  \item \textbf{Log-barrier constraints effectively enforce comfort safety.}
    CSAC-LB reduces the comfort violation rate from $51.6\%$ (RBC) to $4.9\%$
    (final) and achieves zero violations during 24-hour heat waves.  The
    smoothed log-barrier mechanism provides a differentiable safety gradient
    without the oscillation problems of Lagrangian multiplier approaches.

  \item \textbf{CSAC-LB achieves full dominance over the bang-bang baseline.}
    The CSAC-LB final checkpoint improves upon the RBC on \emph{every} key
    performance metric simultaneously: cost ($-5.1\%$), emissions ($-4.8\%$),
    consumption ($-4.9\%$), peak demand ($-5.3\%$), and ramping ($-53.7\%$),
    while reducing comfort violations by $90.5\%$.  This demonstrates that
    the learned policy achieves safety \emph{and} efficiency through
    continuous, modulated cooling---in contrast to the RBC's wasteful
    binary switching.

  \item \textbf{Continuous control outperforms on/off control.}
    The fundamental advantage of the learned policy is its ability to apply
    \emph{graded} cooling actions.  Where the bang-bang RBC oscillates between
    full cooling and no cooling (causing energy waste and grid instability),
    CSAC-LB learns to apply precisely the amount of cooling needed at each
    time step, achieving better comfort with less energy.
\end{enumerate}


% -----------------------------------------------------------------------------
\subsection{Limitations and the Safety--Efficiency Frontier}
\label{sec:temp-limitations}

While CSAC-LB achieves full dominance over the bang-bang RBC baseline in
full-year and 24h heat-wave evaluations, important limitations remain.  The
72h heat-wave scenario reveals that sustained extreme conditions can favour
simple bang-bang control on energy metrics, and the current evaluation uses
only a single RBC baseline.  This section discusses the remaining limitations
and proposes directions for further improvement.

\subsubsection{Barrier Hyperparameters as Safety--Efficiency Control Knobs}

The CSAC-LB framework provides two interpretable parameters---the cost
limit~$d$ and barrier sharpness~$\mu$---that control the position on the
safety--efficiency Pareto frontier.  The current strict configuration
($d = 0.5$, $\mu = 25$) prioritises safety, producing a violation rate of
$4.9\%$ with energy consumption $4.9\%$ below the bang-bang RBC.  Relaxing
these parameters would shift the policy toward lower energy use at the cost
of modestly higher violation rates:

\begin{table}[htbp]
\centering
\caption{Operating points on the safety--efficiency Pareto frontier.  The
  current configuration achieves full dominance over RBC in the full-year
  scenario; relaxed configurations could further reduce energy consumption
  at the cost of modestly higher violation rates.}
\label{tab:pareto-frontier}
\begin{tabular}{lcccl}
\toprule
\textbf{Configuration}
  & $d$ & $\mu$
  & \textbf{Expected Violation}
  & \textbf{Expected Energy} \\
\midrule
Current (strict)     & 0.5  & 25 & $\sim$5\%     & $-$5\% vs bang-bang RBC \\
Mild relaxation      & 2.0  & 15 & $\sim$15--20\% & $-$15--20\% vs bang-bang RBC \\
Moderate relaxation  & 5.0  & 10 & $\sim$20--25\% & further reduction possible \\
\bottomrule
\end{tabular}
\end{table}

The framework thus separates the \emph{mechanism} (barrier function) from
the \emph{target} (how strictly to enforce), giving building operators an
interpretable control interface.

\subsubsection{Proposed Directions for Closing the Gap}
\label{sec:temp-future-directions}

Three principled approaches could further improve results, particularly
for sustained heat-wave scenarios and additional baseline comparisons:

\paragraph{Direction 1: Two-constraint CMDP.}
Formulate energy consumption as a \emph{second} cost constraint alongside
comfort:
\begin{align}
  C_t^{(1)} &= w_{\text{comfort}} \cdot \sum_{i} \mathbf{1}[T_t^{(i)} \notin [20, 26]] \label{eq:cost-comfort} \\
  C_t^{(2)} &= w_{\text{energy}} \cdot \sum_{i} \max\!\left(e_t^{(i)} - e_{\text{RBC},t}^{(i)},\; 0\right) \label{eq:cost-energy}
\end{align}
where $e_{\text{RBC},t}^{(i)}$ is the RBC's energy consumption, serving as
an upper bound on acceptable energy use.  Each constraint receives its own
barrier penalty with independent $d$ and $\mu$ values, enabling simultaneous
safety and efficiency enforcement.  This extends the single-constraint
CSAC-LB to a multi-constraint barrier framework.

\paragraph{Direction 2: Adaptive barrier annealing.}
Anneal the barrier sharpness $\mu$ during training using a curriculum:
\begin{equation}
  \mu(t) = \mu_{\max} \cdot \min\!\left(1,\; \frac{t}{t_{\text{anneal}}}\right)
  + \mu_{\min} \cdot \left(1 - \min\!\left(1,\; \frac{t}{t_{\text{anneal}}}\right)\right)
  \label{eq:barrier-anneal}
\end{equation}
Starting with a relaxed barrier ($\mu_{\min} = 5$) allows the agent to first
learn energy-efficient cooling behaviour; then gradually increasing to
$\mu_{\max} = 25$ tightens the comfort constraint without losing the
efficiency knowledge.  This is analogous to the curriculum induction approach
of Turchetta et~al.~\cite{turchetta2020safe} applied to the barrier
parameter space.

\paragraph{Direction 3: Asymmetric comfort penalties.}
The hot/cold discomfort analysis (Figure~\ref{fig:hot-cold-breakdown})
showed that while CSAC-LB reduces cold discomfort relative to RBC ($0.186$
vs $0.376$), it introduces hot discomfort ($0.025$ vs $0.000$).  Introducing
asymmetric penalties:
\begin{equation}
  C_t^{\text{asym}} = w_{\text{hot}} \cdot \mathbf{1}[T > 26] +
    w_{\text{cold}} \cdot \mathbf{1}[T < 20], \quad
    w_{\text{cold}} \gg w_{\text{hot}}
  \label{eq:asym-cost}
\end{equation}
would discourage the policy from over-cooling, as pushing temperatures below
$20\,^\circ$C would incur a larger penalty than approaching the upper bound
from below.  This directly addresses the mechanism behind the over-cooling
without changing the barrier function itself.

\subsubsection{Summary}

The CSAC-LB final checkpoint achieves full dominance over the bang-bang RBC
baseline in the full-year evaluation, improving on \emph{every} key metric:
$90.5\%$ fewer comfort violations, $5.1\%$ lower cost, $4.8\%$ lower
emissions, $4.9\%$ lower consumption, $5.3\%$ lower peak demand, and $53.7\%$
lower ramping.  Under 24h heat-wave stress testing, CSAC-LB maintains full
dominance with zero violations (vs RBC's $70\%$), $5.0\%$ lower cost, and
$44.6\%$ lower ramping.  The 72h heat-wave scenario reveals a trade-off:
CSAC-LB reduces violations by $83.3\%$ and ramping by $32.6\%$, but at
$\sim$$4.6\%$ higher energy cost.  This demonstrates that safe RL with the log-barrier
mechanism can simultaneously improve both safety and efficiency over a
realistic on/off thermostat baseline, though sustained extreme conditions
remain a challenge.  The three proposed directions---multi-constraint CMDP,
adaptive barrier annealing, and asymmetric penalties---offer principled paths
toward closing the remaining gaps and constitute promising avenues for future
work.


% =============================================================================
% Bibliography entries (to be merged with main bibliography)
% =============================================================================
% @article{zhang2025stems,
%   title={STEMS: Spatial-Temporal Enhanced Safe Multi-Agent Coordination for Building Energy Management},
%   author={Zhang, Huiliang and Wu, Di and Zinflou, Arnaud and Boulet, Benoit},
%   journal={IEEE Internet of Things Journal},
%   year={2025},
%   note={arXiv:2510.14112}
% }
% @article{zhang2024csaclb,
%   title={Constrained Reinforcement Learning for Safe Heat Pump Control},
%   author={Zhang, Baohe and Frison, Lilli and Brox, Thomas and B{\"o}decker, Joschka},
%   journal={arXiv preprint arXiv:2409.19716},
%   year={2024}
% }
% @article{nweye2023citylearn,
%   title={CityLearn v2: Energy-flexible, resilient, occupant-centric, and carbon-aware management of grid-interactive communities},
%   author={Nweye, Kingsley and others},
%   journal={Journal of Building Performance Simulation},
%   year={2024}
% }
% @article{schulman2017ppo,
%   title={Proximal Policy Optimization Algorithms},
%   author={Schulman, John and Wolski, Filip and Dhariwal, Prafulla and Radford, Alec and Klimov, Oleg},
%   journal={arXiv preprint arXiv:1707.06347},
%   year={2017}
% }
% @inproceedings{haarnoja2018sac,
%   title={Soft actor-critic: Off-policy maximum entropy deep reinforcement learning with a stochastic actor},
%   author={Haarnoja, Tuomas and Zhou, Aurick and Abbeel, Pieter and Levine, Sergey},
%   booktitle={International Conference on Machine Learning},
%   year={2018}
% }
% @phdthesis{achiam2021exploration,
%   title={Exploration and Safety in Deep Reinforcement Learning},
%   author={Achiam, Joshua},
%   school={University of California, Berkeley},
%   year={2021}
% }
% @inproceedings{ha2020learning,
%   title={Learning to walk in the real world with minimal human effort},
%   author={Ha, Sehoon and Xu, Peng and Tan, Zhenyu and Levine, Sergey and Tan, Jie},
%   booktitle={Conference on Robot Learning},
%   year={2020}
% }
% @inproceedings{turchetta2020safe,
%   title={Safe Reinforcement Learning via Curriculum Induction},
%   author={Turchetta, Matteo and Kolobov, Andrey and Shah, Shital and Krause, Andreas and Amir, Alekh},
%   booktitle={Advances in Neural Information Processing Systems},
%   year={2020}
% }
% @inproceedings{fujimoto2018td3,
%   title={Addressing Function Approximation Error in Actor-Critic Methods},
%   author={Fujimoto, Scott and van Hoof, Herke and Meger, David},
%   booktitle={International Conference on Machine Learning},
%   year={2018}
% }
